\chapter{Incorpora la retroalimentación}
\label{cap:retro}

\section{Estado de la retroalimentación de la Etapa 1}
\label{sec:estado-retro}

El Informe de Etapa 1 se entregó el 4 de septiembre de 2026 (\href{run:../evidence/informe-etapa-1.pdf}{Informe de Etapa 1 entregado}). La revisión escrita anunciada no llegó como
documento; en la reunión del 5 de octubre el profesor guía informó que la Etapa 1 fue
calificada con 7,0 y que no quedaba retroalimentación pendiente. La única observación abierta
concierne a la implementación: el motor fiel debe alcanzar el rendimiento a un núcleo del
original 2004, revisando sus estructuras de datos (protocolos de la sección
\ref{sec:protocolos}). Su diseño está en la sección
\ref{sec:rendimiento} y en el registro RD-9; la implementación y la medición antes y después
corresponden a la Etapa 3, cuya pauta las pide.

\section{Indicaciones del profesor guía y cambios hechos}
\label{sec:indicaciones}

El cuadro \ref{tab:indicaciones} registra cada indicación del profesor guía desde el inicio
del curso, el cambio que produjo y dónde se ve.

\footnotesize
\begin{longtable}{L{1.6cm}L{3.4cm}L{4.8cm}L{2.4cm}}
\caption{Indicaciones del profesor guía, cambio hecho y dónde se ve.}\label{tab:indicaciones}\\
\toprule
Fecha y medio & Indicación & Cambio hecho & Dónde se ve \\
\midrule
\endfirsthead
\multicolumn{4}{l}{\footnotesize\emph{Cuadro \ref{tab:indicaciones}, continuación}}\\
\toprule
Fecha y medio & Indicación & Cambio hecho & Dónde se ve \\
\midrule
\endhead
\bottomrule
\endfoot
24 ago., mensaje del curso & Límite de treinta páginas; estructura de Cardiff; SWEBOK v4 como
guía de estándares & el cuerpo lleva lo que la pauta evalúa y el detalle va a anexos; cada
vista y decisión cita SWEBOK v4.0a, ISO/IEC/IEEE 42010, IEEE 1016 o UML 2.5.1 & todo el
informe \\
\addlinespace
3 sept., reunión & Derivar los requisitos de lo construido; el umbral RNF01 es del
propietario; no afirmar convergencia & cada decisión cita código o configuración con archivo y
línea; el valor del umbral queda abierto; ciclos fijos, sin prueba de convergencia & cuadro
\ref{tab:decisiones}; sección \ref{sec:experimental}; RD-5 \\
\addlinespace
28 sept., reunión semanal & Mejorar el rendimiento del motor fiel revisando punteros y
estructuras de datos & RNF04 precisado; OE1 con criterio medible y plazos; diseño de la mejora
(D-20, D-21) & sección \ref{sec:rendimiento}; RD-9 \\
\addlinespace
5 oct., reunión & Informe propio por etapa y criterio; diagramas según estándares; decisiones
con su porqué; rendimiento a un núcleo del Fortran & este informe; seis figuras de cuerpo, un cuadro de mecanismos y dos
figuras de anexo con notación declarada; decisiones con alternativa, razón y evidencia; nueve
registros & secciones \ref{sec:decisiones} y \ref{sec:registros} \\
\end{longtable}
\normalsize

\section{Cambios al contenido de la Etapa 1}
\label{sec:cambios-e1}

El Informe de Etapa 1 queda como se entregó. En el documento de trabajo de título
(\href{run:../evidence/documento-tesis-2026-10-04.pdf}{documento de tesis, 4 de octubre de 2026}) su contenido recibió
después cambios que no responden a una observación sobre la Etapa 1, sino a indicaciones
posteriores, a la auditoría de cada cifra contra su fuente y a una revisión adversarial
independiente: el capítulo 2 se completó, porque sus ecuaciones no coincidían con el código; se
agregó RNF09, Mantenibilidad y auditabilidad, que el propietario pide desde el 15 de abril de
2026 y que el diseño no tenía respaldado (cuadro \ref{tab:atributos}); el hito de octubre de
OE1 se acotó y la mejora completa medida pasó a inicios de diciembre, porque no cabe en cuatro
semanas junto con las pruebas de la Etapa 3 (RD-9); se retiraron o corrigieron cifras de los
capítulos 1 y 3, entre ellas la causa del caso de 61 de 168 legisladores; y el registro de
tareas salió del documento hacia la planilla y el repositorio (capítulo \ref{cap:avance}).

\section{Propuesta: el caso chileno como validación en condiciones operativas reales}
\label{sec:propuesta-chile}

Esta propuesta de los autores deriva de la reunión del 5 de octubre y queda pendiente de
confirmación en la revisión del 9 de octubre. En esa reunión el profesor guía describió el
trabajo como dos tareas, mostrar que el programa Fortran se puede replicar y que se puede
mejorar, con una validación cuantitativa contra el Fortran. OE3 enuncia el estudio chileno
«como caso de validación del instrumento», pero lo mantiene como objetivo propio. Se propone
que sea la validación en condiciones operativas reales que pide la pauta de la Etapa 4: si se
confirma, OE3 se reformula o se retira, RF10 a RF12 se trazan a esa validación y la tarea H1
queda resuelta en dos tareas, replicar y mejorar. Mientras no se confirme, los objetivos de la
sección \ref{sec:recap} siguen vigentes.

\section{Decisiones que quedan en manos del profesor guía}
\label{sec:decisiones-guia}

\begin{enumerate}
  \item Confirmación de la propuesta anterior y, con ella, de la descomposición de objetivos
  (tarea H1).
  \item Valor del umbral de fidelidad RNF01, criterio de aceptación de las pruebas
  exhaustivas (tareas C6 y H2); la sección \ref{sec:experimental} propone su forma.
  \item Ratificación del diseño principal de periodización del estudio chileno (D-16).
\end{enumerate}

El 5 de octubre se acordó entregar un informe propio como enlace. El límite de treinta
páginas se aplica por informe y al cuerpo como supuesto de trabajo, no como acuerdo ratificado.
